\documentclass[11pt,a4paper]{article}
\usepackage[T1]{fontenc}
\usepackage[utf8]{inputenc}
\usepackage{lmodern,microtype,amsmath,amssymb,amsthm}
\usepackage[left=26mm,right=26mm,top=23mm,bottom=24mm,headheight=14pt]{geometry}
\usepackage{xcolor,fancyhdr}
\usepackage[hypertexnames=false]{hyperref}
\definecolor{ink}{HTML}{19475A}
\hypersetup{colorlinks=true,linkcolor=ink,citecolor=ink,urlcolor=ink,
pdftitle={A nodal criterion for real zeros of entire functions in a strip},
pdfauthor={Research manuscript prepared with ChatGPT for Keio}}
\pagestyle{fancy}\fancyhf{}\lhead{\small\itshape A nodal criterion for real zeros}
\rhead{\small\thepage}\renewcommand{\headrulewidth}{0.25pt}
\newtheorem{theorem}{Theorem}[section]
\newtheorem{lemma}[theorem]{Lemma}
\newtheorem{proposition}[theorem]{Proposition}
\newtheorem{corollary}[theorem]{Corollary}
\numberwithin{equation}{section}
\renewcommand{\Re}{\operatorname{Re}}\renewcommand{\Im}{\operatorname{Im}}
\setlength{\emergencystretch}{1.5em}
\title{A nodal criterion for real zeros of entire functions in a strip}
\author{Research manuscript prepared with ChatGPT for Keio}
\date{3 November 2026}
\begin{document}
\maketitle
\begin{abstract}
Csordas and Escassut conjectured in 2005 that strict positivity of the first
Laguerre expression for every symmetric imaginary translate of a real entire
function in their strip class forces all its zeros to be real.
We prove this implication without any density or spacing assumption on the
zeros. A nonreal zero forces a simultaneous zero of the real part of a
fixed rotation of the function and its horizontal derivative at a positive
height in the strip, for every rotation with nonzero real component.
The main estimate bounds intersections of this real level with selected
vertical segments by applying Jensen's formula in the complexified height
variable. We also bound the number of nonreal zeros by a sum of local
multiplicities at horizontal tangencies, and characterize when the set of
such tangencies is finite.
\end{abstract}

\section{The shifted Laguerre conjecture}
Let $f$ be a real entire function, so $f(\bar z)=\overline{f(z)}$.
Write $L[g]=(g')^2-gg''$ and
$f_\mu(z)=f(z+i\mu)+f(z-i\mu)$ for real $\mu$. If a real polynomial $p$
has only real zeros $t$, then off its zeros
$L[p](x)/p(x)^2=\sum_t(x-t)^{-2}$, with multiplicity.
The shifts enter when one tries to reverse this implication. For instance
$p(z)=z(z^2+1)$ has $L[p](x)=3x^4+1>0$ for all real $x$, despite
its pair of imaginary zeros \cite{CE}.
For $A\ge0$, let $\mathcal S(A)$ consist of the nonzero real entire functions
with a canonical factorization
\begin{equation}\label{eq:phase-canonical}
 f(z)=Cz^m e^{-az^2+bz}\prod_{\rho} E_1(z/\rho),
 \qquad E_1(w)=(1-w)e^w .
\end{equation}
Here $a\ge0$, $b,C\in\mathbb R$, $C\ne0$, and $m$ is a nonnegative
integer. The product includes every nonzero zero with its multiplicity,
$|\Im\rho|\le A<\infty$, and $\sum_\rho|\rho|^{-2}<\infty$.
Within this class, membership in the Laguerre--P\'olya class
$\mathcal L\text{-}\mathcal P$ means exactly that every zero is real.
The following assertion is Conjecture~3.5 of Csordas and Escassut
\cite[p.~337]{CE}, with its original hypotheses.
\begin{theorem}[Shifted Laguerre converse]\label{thm:original}
If $f\in\mathcal S(A)$ and
\[
             L[f_\mu](x)>0
             \qquad(x\in\mathbb R,\ \mu\in\mathbb R\setminus\{0\}),
\]
then $f\in\mathcal L\text{-}\mathcal P$.
\end{theorem}

The converse question goes back to Csordas, Smith and Varga in 1992.
Their Corollary~I treats functions with finitely many nonreal zeros.
Their Theorem~III allows infinitely many, assuming a positive proportion
of real zeros and at least linear growth of the total zero count in
\emph{each} real direction \cite{CSV}. Csordas and Escassut explicitly
asked whether these density assumptions were necessary
\cite[Remark~3.6]{CE}. As the polynomial above also shows,
strict positivity matters. An exact computation gives
\[
       L[p_\mu](x)=4\{3x^4+(1-3\mu^2)^2\}\ge0
\]
at every point and every shift, with equality at $x=0$,
$\mu=\pm1/\sqrt{3}$ \cite[Remark~3.6]{CE}.

Recent work of Lin Tao established the positive-Gaussian case, the disk
constraint and direction of escape of a level curve, and alignment of
distant nonreal zeros in an escape direction. It also employed a dyadic
selection of scales to control a logarithmic derivative
\cite[Lemma~3.1, Theorem~3.2, Lemma~5.1, Theorem~16.4]{Lin1}.
A September 2026 sequel explicitly retained the full converse as open
\cite[\S11.5, Open E]{Lin2}. We use these geometric ideas and prove the
estimates needed here. The extra step concerns the \emph{number} of
intersections on a vertical segment. A disk estimate for a holomorphic
function of the height rules out the linear number of intersections
that a single escape would force.

\section{Counting intersections on a vertical segment}\label{sec:sparse}
An auxiliary question separates analytic growth from the geometry of the
level curves. Suppose $a=0$ in \eqref{eq:phase-canonical}. For a fixed
angle and a fixed height $A$, can the level
$\Re(e^{i\theta}f)=0$ meet \emph{every} sufficiently far vertical segment
$R+i[0,A]$ at least $\varepsilon R$ times for some $\varepsilon>0$?
For $\cos\theta\ne0$ the answer is negative.

\begin{lemma}\label{lem:sparse}
Let $f$ have the representation \eqref{eq:phase-canonical} with $a=0$.
For every fixed $L>0$ there exist $R_j\to+\infty$ with $f(R_j)\ne0$
such that
\[
 \log\frac{\max_{|h|\le L}|f(R_j+h)|}{|f(R_j)|}=o(R_j).
\]
For every fixed $0<A<L$ and $\cos\theta\ne0$, the number of zeros
of $y\mapsto\Re(e^{i\theta}f(R_j+iy))$ on $[0,A]$, counted with
multiplicity, is $o(R_j)$. The same centers work for every fixed
angle with $\cos\theta\ne0$.
\end{lemma}

\begin{proof}
For fixed
$L>0$, and $R>0$ with $f(R)\ne0$, set
\[
 H_L(R)=\log\frac{\max_{|h|\le L}|f(R+h)|}{|f(R)|}\ge0.
\]
We show that $H_L(R_j)=o(R_j)$ on a sequence $R_j\to\infty$.
Let $N(t)$ count all nonzero zeros with $|\rho|\le t$, including
multiplicities. For $R\in[T,2T]$, $T\ge 2L+1$, the near factors
$|\rho|\le4T$ contribute at most
\[
 \sum_{|\rho|\le4T}
 \left\{\log(1+L/|R-\rho|)+ L/|\rho|\right\}
\]
to the largest increase in logarithmic modulus.
For $|\rho|>4T$, the logarithmic derivative of $E_1(z/\rho)$ is
$z/[\rho(z-\rho)]$. Integrating along the segment from $R$ to $R+h$
bounds each increase by $C_L T/|\rho|^2$. The linear exponential
and monomial contribute $O_{f,L}(1)$. The total far bound is $o(T)$.
Also $\sum_{|\rho|\le4T} 1/|\rho|=o(T)$, by separating finitely many
zeros and bounding the remaining terms by $4T$ times the corresponding
tail of $\sum|\rho|^{-2}$. The integral of
$\log(1+L/|R-\rho|)$ over $[T,2T]$ is $O_L(\log(2T))$ uniformly in
$\rho$, since we may replace $\rho$ by $\Re\rho$ and integrate the decreasing even
function $s\mapsto\log(1+L/|s|)$ over an interval centered at zero
of length $2T$. Thus
\begin{equation}\label{eq:phase-average}
 \frac1T\int_T^{2T} H_L(R)\,dR
 \le o(T)+C_L\frac{N(4T)\log(2T)}{T}.
\end{equation}
The singularities in the upper bound at real zeros on each averaging interval are
integrable. Moreover $\sum_{j\ge1}4^{-j}N(4\cdot2^j)<\infty$. Indeed,
interchanging sums gives each zero of modulus at least one a geometric
tail bounded by $C/|\rho|^2$. Consequently there is a subsequence along
which $j4^{-j}N(4\cdot2^j)\to0$, since $\sum j^{-1}=\infty$.
For $T=2^j$ on this subsequence, \eqref{eq:phase-average} is $o(T)$.
Since $H_L\ge0$, the claimed sequence of centers with $f(R_j)\ne0$
follows.

Now take $0<A<L$. To count level intersections at such a center,
continue the projection in the \emph{vertical} variable by setting
\[
 G_{R,\theta}(w)
 =\tfrac12\{e^{i\theta}f(R+iw)+e^{-i\theta}f(R-iw)\}.
\]
For real $w$ its value is $\Re(e^{i\theta}f(R+iw))$. At the origin it equals
$f(R)\cos\theta\ne0$. Its maximum on $|w|\le L$
is at most $|f(R)|\exp H_L(R)$. Jensen's formula bounds its number
$n_{R,\theta}(A)$ of zeros in the \emph{closed} disk $|w|\le A$, counted
with multiplicities, by
\begin{equation}\label{eq:phase-jensen}
 n_{R,\theta}(A)\log(L/A)
 \le H_L(R)+\log(1/|\cos\theta|).
\end{equation}
Indeed, use circles with radii $r$ increasing to $L$, $A<r<L$,
that avoid zeros. Each zero with $|w|\le A$ contributes at least
$\log(r/A)$. The right side is $o(R_j)$ along the selected centers. These centers were chosen before $\theta$.
\end{proof}

\section{Descent along a level curve}\label{sec:fixedprojection}
Given a real angle $\theta$, put
\[
 u(x,y)=\Re(e^{i\theta}f(x+iy)),\qquad
 v(x,y)=\Im(e^{i\theta}f(x+iy)).
\]
The real entire function whose values on the horizontal line of
height $y$ describe this projection is
\begin{equation}\label{eq:phase-horizontal}
 F_{y,\theta}(z)
 =\tfrac12\{e^{i\theta}f(z+iy)+e^{-i\theta}f(z-iy)\}.
\end{equation}
Thus $F_{y,\theta}(x)=u(x,y)$ for real $x$.

\begin{theorem}\label{thm:phase-main}
Let $f\in\mathcal S(A)$ and fix one real angle $\theta$ with
$\cos\theta\ne0$. Suppose that $u_x(x,y)\ne0$ at every point
where $u(x,y)=0$, $x\in\mathbb R$ and $0<y\le A$. Then all zeros
of $f$ are real. Equivalently, if $f$ has a nonreal zero, then for
every such fixed angle there exist $x\in\mathbb R$ and $0<y\le A$
with $F_{y,\theta}(x)=F_{y,\theta}'(x)=0$.
\end{theorem}

When $f$ has a zero, no $F_{y,\theta}$ with $y>0$ can vanish
identically. Otherwise \eqref{eq:phase-horizontal} yields
$f(z+2iy)=-e^{-2i\theta}f(z)$ for all $z$, so vertical translation
repeats a zero outside the strip. The point detected in
Theorem~\ref{thm:phase-main} is therefore a real zero of multiplicity
at least two of a nonzero entire function.

\subsection{Descent from a nonreal zero}
\noindent\textit{Proof of Theorem~\ref{thm:phase-main}.} Assume the transversality hypothesis in Theorem~\ref{thm:phase-main}.
By the reality of $f$, the presence of a nonreal zero implies the
presence of an upper-half-plane zero. Let $x_0+iy_0$ be any such zero.
The case $A=0$ is vacuous, so suppose $A>0$. Locally the level set
$u=0$ through $(x_0,y_0)$ is a graph $x=X(y)$ by the implicit
function theorem. Continue this graph downward on its maximal interval
$(y_*,y_0]$, $y_*\ge0$. The simple continuation used here depends
only on $u_x\ne0$ on this level in $0<y\le A$. The
Cauchy--Riemann equations give
\begin{equation}\label{eq:phase-descent}
 X'(y)=-u_y/u_x,\qquad
 \frac{d}{dy}v(X(y),y)=\frac{u_x^2+u_y^2}{u_x}.
\end{equation}
The denominator has fixed nonzero sign on the graph, and $v(X(y_0),y_0)=0$.
Hence $v u_x<0$ below the starting zero and $|v|$ strictly increases
as the ordinate decreases. In particular $|v|$ stays bounded away
from zero near the lower end of the graph.

There can be no finite accumulation point at that end. If $y_*=0$
and $X(y_n)\to x_*$ for some $y_n\downarrow0$, continuity gives
$0=u(x_*,0)=\cos\theta\,f(x_*)$ and thus $v(x_*,0)=0$, a
contradiction. If $y_*>0$, the implicit function theorem gives a
local graph through any finite accumulation point at height $y_*$.
It meets the original graph near that point and continues it below
$y_*$ by local uniqueness. Consequently $X(y)\to+\infty$ or
$X(y)\to-\infty$ as $y\downarrow y_*$. The sign follows from
continuity once $|X(y)|\to\infty$.

Graphs descended from distinct upper zeros are disjoint. At an
intersection uniqueness makes the graphs coincide for as long as
both are defined. In particular the union on the ordinate interval
between their starting zeros lies on one regular graph. The
strictly monotone function $y\mapsto v(X(y),y)$ cannot vanish at
both ends. If the starting heights agree, uniqueness instead gives
the same starting zero.

At any point on a descended graph below its starting zero, the
partial derivative at fixed $x$ satisfies
\begin{equation}\label{eq:phase-negative}
 \partial_y\log|f(x+iy)|=\frac{u u_y+v v_y}{u^2+v^2}
 =\frac{u_x}{v}<0 .
\end{equation}

\subsection{How many branches must cross a distant segment?}
Logarithmic differentiation of \eqref{eq:phase-canonical} gives, in
$y>0$ and away from the zeros, a contribution $2ay\ge0$ to
$\partial_y\log|f|$ from the Gaussian and a contribution
$y/((x-t)^2+y^2)>0$ from each real zero $t$, including $t=0$.
The linear exponential has zero contribution. A conjugate pair
$t\pm i\eta$, $\eta>0$, contributes
\begin{equation}\label{eq:phase-disk}
 \frac{2y((x-t)^2+y^2-\eta^2)}
 {((x-t)^2+(y-\eta)^2)((x-t)^2+(y+\eta)^2)} .
\end{equation}
The convergence factors in each conjugate pair have real linear
coefficient and hence zero contribution. The differentiated canonical
product converges absolutely on every compact set away from its zeros.
Equations \eqref{eq:phase-negative} and \eqref{eq:phase-disk} show
that every point $(x,y)$ of a descended graph below its starting zero
satisfies
\begin{equation}\label{eq:phase-cover}
 (x-t)^2+y^2<\eta^2\le A^2
\end{equation}
for some upper-half-plane zero $t+i\eta$.

We may take the escape to be rightward. For a left escape, replace $f(z)$ by
$f(-z)$ and $\theta$ by $-\theta$. The new real part at $(x,y)$ is
$u(-x,y)$, and the reflected function still belongs to $\mathcal S(A)$. The original escaping graph meets every vertical line to
the right of its starting point. By \eqref{eq:phase-cover}, every real
number beyond some fixed $X_0$ lies within $A$ of the real coordinate
of an upper zero. Choose $X_1>X_0$ with $f(X_1)\ne0$.
For $R>X_1+2A$, all intervals $[t-A,t+A]$ obtained from upper zeros
with $X_1\le t\le R$ together cover $[X_1+A,R-A]$. Their total
length gives at least $(R-X_1-2A)/(2A)$ distinct upper zeros in that
vertical slab.

Every left-escaping graph from one of these zeros meets the compact
segment $X_1+i[0,A]$. The intersections from different zeros must be
distinct. They are zeros of $y\mapsto u(X_1,y)$, which is real
analytic on a neighborhood of $[0,A]$ and nonzero at zero because
$u(X_1,0)=\cos\theta\,f(X_1)$. There are at most a fixed finite
number of them. All remaining graphs from the slab meet the segment
$R+i[0,A]$ at distinct points. In particular the assumption of one
escape implies, for every sufficiently large $R$,
\begin{equation}\label{eq:phase-many}
 \#\{y\in(0,A]:u(R,y)=0\}\ge R/(2A)-O_{f,A,\theta}(1).
\end{equation}


When $a=0$, Lemma~\ref{lem:sparse} supplies a sequence of these
segments with only $o(R)$ intersections. This contradicts
\eqref{eq:phase-many}.

For completeness consider $a>0$. The product including $Cz^m e^{bz}$
but omitting the Gaussian has logarithmic maximum $o(r^2)$.
The global bound $\log|E_1(w)|\le C_0|w|^2$ follows by expansion at zero
and $\log(1+|w|)+|w|$ elsewhere. A finite product together with
$Cz^m e^{bz}$ has logarithmic maximum $O(r+\log(2+r))$, while each
remaining tail is bounded by $C_0r^2\sum_{\rho\text{ in tail}}|\rho|^{-2}$.
It follows that $f(x+iy)\to0$ uniformly for $0\le y\le A$ as
$|x|\to\infty$. This directly contradicts the positive lower bound
on $|v|$ at the escaping end of a descended graph, and completes
the proof of Theorem~\ref{thm:phase-main}.\hfill$\square$

\begin{proof}[Proof of Theorem~\ref{thm:original}]
Take $\theta=0$. Since $f_\mu(x)=2F_{\mu,0}(x)$, a common real zero of
$F_{\mu,0}$ and $F'_{\mu,0}$ makes $L[f_\mu](x)=0$ there. Strict
positivity in the original conjecture excludes every such point with
$0<\mu\le A$. Theorem~\ref{thm:phase-main} proves the claim.
\end{proof}

If $f$ has at least one zero and all its zeros are real, the converse
to the geometric criterion holds as well. In $y>0$ the product gives
$\partial_y\log|f|=2ay+\sum_t y/((x-t)^2+y^2)>0$. At a point with $u=0$
the function $f$ is nonzero and $u_x/v=\partial_y\log|f|>0$.
Thus no positive-height projection has a simultaneous zero with its
horizontal derivative.

\section{Counting upper zeros by tangencies}
\label{sec:counting}
Fix a real $\theta$ with $\cos\theta\ne0$ and write
$F=e^{i\theta}f$, $u=\Re F$ and $v=\Im F$. For $f\in\mathcal S(A)$,
let
\[
 B=\{b=x+iy:0<y\le A,\ u(b)=u_x(b)=0\},\qquad
 m_b=\operatorname{ord}_b(f-f(b)).
\]
Thus $m_b=1$ at a tangency where $f'(b)\ne0$. Constant functions have
no zeros or tangencies under the stated phase condition and may be
omitted. Let $N_+(I)$ count zeros of $f$ with multiplicity in
$\{x+iy:x\in I,\ 0<y\le A\}$. For a real $s$, put
\[
 G_s(w)=\tfrac12\big(e^{i\theta}f(s+iw)+e^{-i\theta}f(s-iw)\big).
\]
If $G_s\not\equiv0$, let $V_s$ count its zeros with multiplicity on
the real segment $(0,A]$.

\begin{theorem}\label{th:finite-tangencies}
If $c<d$, both wall functions $G_c,G_d$ are nonzero as entire
functions, and the sum on the right is finite, then
\begin{equation}\label{eq:bounded-window}
 N_+((c,d))\le V_c+V_d+
 \sum_{\substack{b\in B\\c<\Re b<d}}m_b .
\end{equation}
If $B$ is finite, then
\begin{equation}\label{eq:global-tangencies}
 N_+(\mathbb R)\le\sum_{b\in B}m_b .
\end{equation}
\end{theorem}

\begin{proof}
A zero in $B$ is counted with its own multiplicity $m_b$. At each
remaining upper zero, $u_x\ne0$, so the zero is simple. Continue its
nodal graph $x=X(y)$ downwards while $u_x\ne0$. Since
$d v(X(y),y)/dy=(u_x^2+u_y^2)/u_x$, the sign of $v$ below the
initial zero is fixed and $|v|$ is strictly increasing during descent.
The interiors of the graphs for distinct zeros cannot meet. Local
uniqueness would identify them back to the lower of their starting
heights, while $v$ cannot vanish twice on such a graph.

First stop the graph on contact with either wall $x=c,d$. A graph
confined to the closed rectangle has an accumulation point $p$ at its
lower endpoint. Strict monotonicity and continuity give $v\to v(p)\ne0$.
In fact the graph converges to $p$. For any sufficiently small disk
centered at $p$, the minimum of $|F-F(p)|$ on its boundary is positive.
Once the graph enters that disk sufficiently late it cannot leave,
since its values approach $F(p)$. The limit cannot have ordinate zero.
There $u(p)=\cos\theta\,f(p)=0$ would give $v(p)=0$. Away from a wall,
positive ordinate and $u_x(p)\ne0$ allow continuation by the implicit
function theorem. Thus $p$ lies on a wall or belongs to $B$.

At any such endpoint $p$, write $v_p=v(p)$ and
$m_p=\operatorname{ord}_p(f-f(p))$. In a local conformal coordinate
$\zeta$ centered at $p$ one has $F-F(p)=\zeta^{m_p}$. For the terminal
portion of an incoming graph one has
\[
 F-F(p)=-i\,\operatorname{sgn}(v_p)t,\qquad t\downarrow0.
\]
This half-ray has exactly $m_p$ local preimage arcs. Distinct starting
zeros have disjoint graph interiors and therefore use distinct terminal
arcs. No graph can end at a zero of $f$ since $v_p\ne0$. An interior
point of $B$ can consequently account for at most $m_p$ zeros, including
its own multiplicity if it is a zero. At a wall $\Re p=s$, the Taylor
series of $F(z)-F(p)$ starts in degree $m_p$, and $u(p)=0$. Thus
$G_s$ has a zero of order at least $m_p$ at $\Im p$, since on the real
axis $G_s(y)=u(s,y)$. This counts every incident arc even at a wall
tangency. Zeros on the walls are excluded from the left side of
\eqref{eq:bounded-window}, so all endpoints are accounted for.

Now assume $B$ is finite and continue each good graph without walls.
If it has a finite accumulation point at its lower endpoint, the same
argument gives a limit in $B$. Otherwise $X(y)$ tends to $+\infty$ or
$-\infty$. At every point below the initial zero,
$\partial_y\log|f|=u_x/v<0$, with the real coordinate held fixed.
The conjugate-pair calculation in Section~\ref{sec:fixedprojection} applies
unchanged. If a graph escapes right, each sufficiently large $x$ is
therefore within distance $A$ of the real coordinate of an upper zero.
Choose $c$ fixed and sufficiently large with $f(c)\ne0$. For $d>c+2A$,
the length of $(c+A,d-A)$ then gives
$N_+((c,d))\ge(d-c-2A)/(2A)$, with $A>0$ understood. For $a=0$ the
centers of Lemma~\ref{lem:sparse} have $f(d)\ne0$ and $H_L(d)=o(d)$.
Since $G_d(0)=\cos\theta\,f(d)$ and
$\max_{|w|\le L}|G_d(w)|\le|f(d)|e^{H_L(d)}$, Jensen's formula gives
$V_d=o(d)$ along these centers. This contradicts
\eqref{eq:bounded-window} because the sum over all of $B$ is finite.
The left-escape argument uses the same estimates at negative centers,
obtained by applying the center estimate to $f(-z)$. For $a>0$,
uniform decay on the strip contradicts the positive lower bound for
$|v|$ along an escaping graph. Hence every good graph terminates in
$B$, and the local count at its endpoint proves
\eqref{eq:global-tangencies}. The case $A=0$ is immediate.
\end{proof}

More precisely, under the finite-$B$ hypothesis, for each $0<h\le A$
the upper zeros of height at most $h$, counted with multiplicity,
are bounded in number by $\sum_{b\in B,\,\Im b\le h}m_b$.
Indeed every graph is charged at a lower height than its initial zero,
and a zero in $B$ is counted at its own height.

Equality occurs already for $f(z)=z^2+1$, $A=1$, $\theta=0$. Here
$u=x^2-y^2+1$, $u_x=2x$, and $B=\{i\}$ with $m_i=1$. The same
polynomial shows why the phase restriction matters. At $\theta=\pi/2$
one has $u=-2xy$ and $u_x=-2y$, so $B$ is empty despite the upper zero.
\begin{corollary}\label{cor:finite-iff}
Let $f\in\mathcal S(A)$ have at least one zero, and fix
$\cos\theta\ne0$. Then $f$ has finitely many nonreal zeros if and
only if $B$ is finite. In that case the total number of nonreal zeros,
with multiplicity, is at most $2\sum_{b\in B}m_b$.
\end{corollary}
\begin{proof}
Theorem~\ref{th:finite-tangencies} gives one implication and the count
by conjugation. For the converse, suppose the nonreal zeros form a
finite set. Choose $M$ large enough that every interval
$[\Re\rho-A,\Re\rho+A]$ centered at a nonreal zero lies in $[-M,M]$.
For $|x|>M$ and $y>0$, each real zero and each conjugate pair makes a
strictly positive contribution to $\partial_y\log|f(x+iy)|$, and the
Gaussian contribution is nonnegative. At least one zero is present.
Thus $\partial_y\log|f|>0$ there. At a nodal point outside the
rectangle one has $v\ne0$ and $\partial_y\log|f|=u_x/v$, which excludes
$u_x=0$. Consequently $B\subset[-M,M]+i(0,A]$.

Suppose $B$ has an accumulation point $p$ in the closed rectangle,
possibly on its real edge or top edge. Then $u(p)=0$. The local
coordinate $F-F(p)=\zeta^m$ used in the proof of
Theorem~\ref{th:finite-tangencies} parametrizes $u=0$ near $p$ by
finitely many regular real-analytic arcs through $p$. Shrinking the
neighborhood, $F'$ has no zero there except possibly $p$. At any
other point of these arcs, $u_x=0$ precisely when the height on the
arc has derivative zero. On an arc containing infinitely many of
the accumulating points, the height derivative therefore vanishes
identically by real analyticity through $p$. Thus $u$ vanishes on a
horizontal interval at a height $\eta\in[0,A]$. Analytic continuation
in $z$ and the reality of $f$ give
\[
 f(z+2i\eta)=-e^{-2i\theta}f(z)\qquad(z\in\mathbb C).
\]
For $\eta=0$ this forces $f\equiv0$ because $\cos\theta\ne0$.
For $\eta>0$ it propagates any zero by arbitrary integer multiples
of $2i\eta$, contradicting the fixed strip. Hence $B$ has no
accumulation point in the compact rectangle and is finite.
\end{proof}

The zero-free exception is explicit. Then
$f(z)=C\exp(-az^2+bz)$ with $a\ge0$. For $a>0$ the identity
$\partial_y\log|f|=2ay$ gives $B=\varnothing$. For $a=0$ the set $B$
is the union of entire horizontal lines at heights $0<y\le A$ for
which $\cos(\theta+by)=0$.

\section{Which operators detect nonreal zeros?}
\subsection{A family of detecting differential operators}

In Open Problem~3.3 of their 2005 paper, Csordas and Escassut ask which
Laguerre--P\'olya symbols $\varphi$ of order less than two have the
following property. For every $f\in\mathcal S(A)$, the strict inequalities
$L[\varphi(\mu D)f](x)>0$ for all real $x$ and all real $\mu\ne0$
imply $f\in\mathcal L\text{-}\mathcal P$. Here $D=d/dx$.
The strip argument gives an exact answer within one family.

\begin{proposition}
Let $C,b,c,\theta$ be real with $Cc\ne0$. The order-one symbol
$\varphi(z)=C e^{bz}\cos(cz+\theta)$ has the stated property
if and only if $\cos\theta\ne0$.
\end{proposition}

\begin{proof}
Put $U+iV=e^{i\theta}f$ and assume first that $\cos\theta\ne0$.
For $y=c\mu$ the entire-function translation identity gives
\[
 [\varphi(\mu D)f](x)= C\,U(x+b\mu,y)
 \qquad(x,\mu\in\mathbb R).
\]
Since $c\ne0$, every positive height occurs, and translation of the
real argument is bijective at each fixed height. Thus the strict
inequalities exclude a common zero of $U,U_x$ at any positive
height. Theorem~\ref{thm:phase-main} gives the conclusion.

If $\cos\theta=0$, choose $f(z)=z^2+1$, which belongs to $\mathcal S(1)$.
Then $[\varphi(\mu D)f](x)=-2C c\mu\sin\theta\,(x+b\mu)$
has Laguerre expression $4C^2c^2\mu^2>0$ for every $\mu\ne0$,
although the two zeros of $f$ are nonreal.
\end{proof}

All these symbols are in the Laguerre--P\'olya class by the canonical
product for cosine and closure under multiplication by real linear
exponentials. The restriction $Cc\ne0$ excludes the zero symbol and
pure translations. For a nonzero pure translation symbol
$K e^{bz}$ the property fails as well. For $p(z)=z(z^2+1)$ one has $L[p](x)=3x^4+1>0$, and a real
translation and multiplication by a nonzero real scalar preserve
strict positivity. For the zero symbol the strict premise cannot hold.

The obstruction at $\cos\theta=0$ extends to every nonzero symbol with
$\varphi(0)=0$. Write $\varphi(z)=\sum_{j\ge k}a_jz^j$ with
$k\ge1$ and $a_k\ne0$. The polynomial $p(z)=z^{k-1}(z^2+1)$ gives
\[
 \varphi(\mu D)p(x)=(k+1)!\mu^k(a_kx+a_{k+1}\mu).
\]
Its Laguerre expression is strictly positive at every $x$ when
$\mu\ne0$, while $p(i)=p(-i)=0$. In particular, $\varphi(0)\ne0$
is necessary for any nonzero symbol in Open Problem~3.3.
\subsection{Why a dense set of heights does not suffice}
Fix $A>0$ and any prescribed countable set $D\subset(0,\infty)$.
Choose $a\in(0,A]$ with $a/\sqrt3\notin D$. Then $p_a(z)=z(z^2+a^2)$
has nonreal zeros in the strip, but
\[
 L[(p_a)_\mu](x)=4\{3x^4+(a^2-3\mu^2)^2\}>0
 \qquad (x\in\mathbb R,\ |\mu|\in D).
\]
Thus even strict tests at all rational heights can miss a nonreal pair.
For the simpler nodal test, each height $0<h\le A$ can be uniquely
necessary. At phase zero, $f(z)=z^2+h^2$ has
$u=u_x=0$ in the positive half-strip precisely at $(0,h)$.

\subsection{The exceptional phase under positive zero-density assumptions}
The phase restriction can already be seen in the printed rotating
variants of the 1992 results. Corollary~2.6 and Theorem~3.5 of
\cite[pp.~386, 394--395]{CSV} permit an arbitrary fixed phase in
$f_{\mu,\theta}=2F_{\mu,\theta}$. For the exceptional phase
$\theta=\pi/2$ the polynomial $z^2+1$ gives
$f_{\mu,\pi/2}(x)=-4\mu x$ and $L[f_{\mu,\pi/2}](x)=16\mu^2>0$.
There is also an example meeting the bilateral density conditions of
their Theorem~3.5. Put
\[
 q(z)=\cos(2z)+\cos z-1.
\]
Its zeros solve $\cos z=(-1\pm\sqrt{17})/4$. One value is in $(0,1)$
and the other is below $-1$. Each half-open period of length $2\pi$
in the real coordinate thus contains two real zeros and a conjugate
pair at fixed nonzero heights, all simple. The fraction of real zeros
tends to $1/2$ in both directions and the total count per unit length
tends to $2/\pi$. The real even function $q$ has order one and its
Hadamard factorization places it in $\mathcal S(A)$ for a finite $A$.
Yet for every $\mu\ne0$,
\[
 2F_{\mu,\pi/2}(x)=2\sin x\,[2\sinh(2\mu)\cos x+\sinh\mu].
\]
Here the projection is taken for $f=q$. Its zeros are given by
$\sin x=0$ or $\cos x=-1/(4\cosh\mu)$. All of them are real and
simple, since the latter value lies strictly between $-1$ and $1$.
For this projected function the order-one canonical product gives
strict positivity of its Laguerre expression off its zeros by the
reciprocal-square sum, and at its zeros by the nonzero derivative.
This identifies an exception in the two phase-dependent assertions.
The zero phase in Theorem~\ref{thm:original} is covered by
Theorem~\ref{thm:phase-main}.

\paragraph{Relation to the source corpus.}
Manuscript~004 of the supplied corpus studies level-point counts in
compact windows for prime Dirichlet series and uses Jensen averages to
study their zero distributions \cite{corpus}. The corresponding object
here is the number of intersections with a segment. Complexifying the
height lets a disk estimate count them directly. The canonical-product
and geometric arguments above are self-contained.

\begin{thebibliography}{9}\raggedright
\bibitem{CSV} G.~Csordas, W.~Smith and R.~S.~Varga,
Level sets of real entire functions and the Laguerre inequalities,
\emph{Analysis} \textbf{12} (1992), 377--402.
\href{https://doi.org/10.1524/anly.1992.12.34.377}{doi:10.1524/anly.1992.12.34.377}.
\bibitem{CE} G.~Csordas and A.~Escassut,
The Laguerre inequality and the distribution of zeros of entire functions,
\emph{Ann. Math. Blaise Pascal} \textbf{12} (2005), 331--345.
\href{https://doi.org/10.5802/ambp.210}{doi:10.5802/ambp.210}.
\bibitem{Lin1} Lin Tao,
\emph{Partial Resolutions of the Shifted Laguerre Converse: Gaussian Drift and
Strict Directional Zero-Density Barriers}, revision 1.3, 2 September 2026.
\href{https://doi.org/10.5281/zenodo.22242139}{doi:10.5281/zenodo.22242139}.
\bibitem{Lin2} Lin Tao,
\emph{Collective Screening and Real-Mesh Obstructions in the Shifted Laguerre
Converse}, self-containment revision, 15 September 2026.
\href{https://doi.org/10.5281/zenodo.22763987}{doi:10.5281/zenodo.22763987}.
\bibitem{corpus} QEEDDD, \emph{Verifiable Mathematical Reasoning Corpus},
manuscript 004, \emph{Prime-Zeta Divisor Tomography: Local Zero Processes,
Frozen-Edge Rigidity, and Derivative Caustics}, revision
\href{https://github.com/QEEDDD/Verifiable-Mathematical-Reasoning-Corpus/tree/ecd87de3348dcb706a76bb12fde722d7a348a012}{\texttt{ecd87de3348dc}}.
\end{thebibliography}
\end{document}
